\documentclass[11pt,letterpaper,sans]{moderncv}

% moderncv styles: casual, classic, banking, oldstyle, fancy
% Order matters: the style file snapshots the active theme colour with
% \colorlet, so \moderncvcolor must be set first or the section headings
% and rules silently fall back to black.
\moderncvcolor{blue}
\moderncvstyle{classic}

% Page geometry
\usepackage[scale=0.9]{geometry}

% Encoding / typography
% lmodern must come before fontenc: it swaps bitmap Computer Modern for the
% scalable Type 1 Latin Modern faces, which microtype's font expansion
% requires and which keep the PDF crisp and text-selectable.
\usepackage{lmodern}
\usepackage[T1]{fontenc}
\usepackage{ragged2e}
\usepackage{etoolbox}
\usepackage{needspace}

% moderncv typesets \cventry's description inside a tabular and a minipage,
% neither of which can split across a page boundary. A long entry therefore
% jumps wholesale to the next page and leaves a gap behind it. \cvjob emits
% only the header row and \cvbullet emits each bullet as its own row (moderncv's
% own \cvlistitem), so LaTeX can break between bullets instead.
%
% \cvitem ends each row with \par\addvspace, which leaves a legal breakpoint
% directly after a header, so \nopagebreak cannot stop a header being stranded
% at the foot of a page: the break is taken at the glue before the penalty.
% \Needspace* reserves room for the header plus a few lines of bullet and moves
% the whole group to the next page when that does not fit. It must be the
% starred form -- moderncv sets \raggedbottom (moderncv.cls), and plain
% \needspace's stretchy-glue trick misfires there, breaking the page early and
% leaving a large gap regardless of how much room actually remains.
\newcommand{\cvjob}[4]{%
  \Needspace*{5\baselineskip}%
  \cventry{#1}{#2}{#3}{#4}{}{}}
% \small matches the size \cventry applies to its description block, and
% \RaggedRight keeps the narrow main column from developing rivers of
% whitespace the way full justification does.
\newcommand{\cvbullet}[1]{\cvlistitem[.1em]{\small\RaggedRight #1}}

% Reduce blocky justification in bullet lists
\AtBeginEnvironment{itemize}{\RaggedRight}
\hyphenpenalty=7000
\exhyphenpenalty=7000
\righthyphenmin=4
\emergencystretch=2em

% Personal data
\name{Peter}{Selby}
\title{Software Engineer}
\address{Newcastle, WA}{}{}
\phone[mobile]{(805)~604~5298}
\email{peter@selby.ca}

\begin{document}
\makecvtitle
\pagestyle{empty}

\section{Summary}
\cvitem{}{Software engineer with over a decade of experience designing and building large-scale distributed systems at Google, ByteDance, Apple, and Amazon. Built the ingestion pipelines and reporting platform behind Amazon's shipment-level retail P\&L, and the automated rollout systems that safely ship real-time communication services to 100M+ users at Google. Strongest in backend and data-intensive systems.}

\section{Technical Skills}
\cvitem{Languages}{Go, Scala, Python, Java, C++, Clojure, Lisp, SQL, JavaScript}
\cvitem{Backend \& Systems}{Service and API design, microservices, gRPC, Protocol Buffers, REST, concurrency, profiling, encryption, Linux \& Unix, NixOS}
\cvitem{Data}{Data warehouse design, distributed ETL, Hadoop, Spark, Amazon Redshift, PostgreSQL, MySQL, Oracle SQL; schema design, query optimization, benchmarking}
\cvitem{Cloud \& Infrastructure}{Kubernetes, Docker, AWS and comparable internal cloud platforms at Google and ByteDance; S3 and object storage, Lambda and equivalent serverless runtimes, DynamoDB, Redis, managed PostgreSQL}
\cvitem{Delivery \& Operations}{CI/CD pipelines, ArgoCD, Flux, GitOps, progressive rollouts and canarying, Prometheus, Grafana, production oncall and incident response}

\section{Experience}

\cventry
  {2025--Present}
  {Career Break}
  {}
  {}
  {}
  {
    Primary caregiver for my son while my wife supported her father through a
    serious illness overseas. Actively seeking my next role.
  }

\cvjob{2023--2025}{Site Reliability Engineer}{ByteDance Inc}{Bellevue, WA}
\cvbullet{Operated and evolved real-time communication (RTC) infrastructure for
  large-scale live streaming, running on Kubernetes across AWS and ByteDance's
  internal cloud platform.}
\cvbullet{Designed and implemented a distributed rolling upgrade system for
  worldwide RTC media ingress nodes, enabling safe and automated fleet upgrades
  with minimal redundancy or capacity loss. Reduced deployment cycle from
  multiple weeks to a 24-hour release, increasing development velocity and
  service reliability.}
\cvbullet{Built and operated cloud services in Go and Python, using ArgoCD and
  Flux for GitOps-based continuous deployment. Worked directly with managed
  cloud primitives: object storage, serverless functions, Redis, and managed
  PostgreSQL.}
\cvbullet{Carried production oncall for large-scale distributed services here
  and in previous roles, diagnosing and resolving live failures with Prometheus
  and Grafana---debugging experience that informs how I design and instrument
  the systems I build.}

\cventry
  {2021--2023}
  {Career Break}
  {}
  {}
  {}
  {
    Extended travel and time with family, after ten years of continuous work.
  }

\cvjob{2016--2021}{Site Reliability Engineer}{Google Inc}{Kirkland, WA}
\cvbullet{Evaluated Google's several internal configuration management
  platforms and selected the one to build on, then rewrote multiple RTC
  services---telephony, chat, and video---to safely consume, validate, apply,
  and monitor its live configuration updates, with support for canaries, A/B
  testing, and automated rollbacks.}
\cvbullet{Partnered with 5+ development teams to standardize that configuration
  logic and ensure consistent behavior across heterogeneous systems serving
  100M+ users.}
\cvbullet{Redesigned and replaced the cloud rollout system for 20+ RTC services,
  implementing multi-stage canarying, key-metric monitoring, and automated
  regression analysis, enabling data-driven rollout and rollback decisions.
  Reduced rollout risk and improved release velocity across a multi-team
  production environment.}
\cvbullet{Built and ran services on Google's internal container orchestration
  and cloud platform---the direct ancestor of Kubernetes---working with its
  equivalents of managed storage, queueing, and database offerings, service
  authentication, and capacity planning.}
\cvbullet{Collaborated with cross-functional teams to refactor core RTC services
  for performance, scalability, and observability, working extensively with gRPC
  for service communication.}

\cvjob{2015--2016}{Full Stack Engineer}{Apple Inc}{Cupertino, CA (Remote)}
\cvbullet{Built internal web applications supporting Apple Maps Transit,
  including an infinite-scrolling POI data editor and a route metadata editor
  for ingested transit routes, accelerating data-entry workflows ahead of the
  Apple Transit launch.}
\cvbullet{Designed features end to end, from PostgreSQL schema through Scala
  services to the JavaScript UI.}
\cvbullet{Designed data models and validation workflows enabling rapid editing
  and publishing of curated transit data produced by on-site contractors.
  Optimized queries, schemas, and batch jobs for complex transit data
  relationships.}

\cvjob{2011--2015}{Software Development Engineer}{Amazon.com}{Seattle, WA}
\cvbullet{Built the ingestion pipelines behind Amazon's shipment-level sales
  data: distributed ETL in Hadoop and SQL drawing from Oracle transactional
  systems and producing a single per-shipment table consumed across the
  business.}
\cvbullet{Designed and implemented the allocation and reconciliation algorithms
  at its core, attributing diverse cost and revenue components---liquidations,
  returns, shipping, and promotions---to individual shipments. Partnered with
  finance and operations stakeholders to define heuristics and validate
  financial correctness.}
\cvbullet{Extended data coverage to new geographies and business domains,
  including FBA and 3P, and launched the data-processing pipeline for Amazon
  China.}
\cvbullet{Joined a new team as its second engineer to turn that data into usable
  reporting: retail P\&L aggregated at multiple levels and reconciled against
  accounting systems, delivered through a drill-down report portal.}
\cvbullet{Built an interactive query service in Scala that opened shipment-level
  data to a far wider audience---until then it was reachable only by writing SQL
  against the warehouse directly.}
\cvbullet{Helped bootstrap that team's roadmap, architecture, and
  technology-stack decisions, and launched the platform as an AWS cloud-native
  service on Redshift and DynamoDB---among the first at Amazon to adopt
  Redshift.}

\section{Education}

\cventry
  {2008--2011}
  {Master of Science, Computer Science}
  {University of British Columbia}
  {Vancouver, BC}
  {}
  {
    Focus areas: distributed systems, operating systems, and cloud computing.
    Selected projects: two variants of a peer-to-peer distributed filesystem
    (directory-server and DHT-based) exploring metadata scalability, and a
    proof-of-concept alternative to object-oriented programming in Lisp, in
    which objects emerged dynamically via pattern matching.
  }

\cventry
  {2003--2008}
  {Bachelor of Science, Honours, Computer Science}
  {University of Winnipeg}
  {Winnipeg, MB}
  {}
  {}

\end{document}